\documentclass[../DoAn.tex]{subfiles}

\begin{document}

\section{Kết luận chung}
\label{section:general-conclusion}

Đồ án tập trung vào bài toán trích xuất, truy xuất và hỏi đáp thông tin trên tài
liệu PDF theo hướng câu trả lời có căn cứ từ tài liệu nguồn. Hệ thống được tổ
chức thành hai luồng chính: luồng tiếp nhận và trích xuất PDF, và luồng
xử lý truy vấn trên dữ liệu đã trích xuất. Cách tách này giúp làm rõ vai trò của
từng tầng: tài liệu PDF trước hết phải được chuyển thành block, bảng, ô và
siêu dữ liệu nguồn; sau đó các đơn vị này mới được dùng để lập chỉ mục, truy xuất
bằng chứng, sinh câu trả lời và tạo dẫn chứng.

Kết quả thực nghiệm cho thấy chất lượng trích xuất PDF ảnh hưởng trực tiếp đến
chất lượng hỏi đáp. Định tuyến theo vùng nội dung giúp hệ thống ghi nhận rõ hơn
các vùng bảng và vết xử lý; Hybrid TATR cải thiện việc tái dựng cấu trúc và gán
văn bản vào ô; chia đoạn có nhận thức bảng giúp hệ thống truy xuất được bằng
chứng ở mức hàng và ô thay vì chỉ dừng ở mức trang hoặc đoạn văn. Những kết
quả này củng cố nhận định rằng trích xuất PDF không phải là bước tiền xử lý phụ,
mà là nền tảng quyết định khả năng truy xuất và kiểm chứng ở các tầng sau.

Ở luồng truy vấn, kết quả cho thấy việc kết hợp tín hiệu từ khóa, tín hiệu ngữ nghĩa
và siêu dữ liệu nguồn phù hợp với tài liệu quy định và tài liệu bán hình thức. Cấu hình
đầy đủ đạt kết quả khả quan trên tập đánh giá chính QCDT, đặc biệt ở
khả năng tìm bằng chứng và tạo dẫn chứng. Tuy nhiên, kết quả trên QASPER cũng
cho thấy hệ thống chưa đủ mạnh cho các tài liệu khoa học dài và các câu hỏi cần
tổng hợp nhiều bằng chứng. Vì vậy, đóng góp của đồ án nên được hiểu trong phạm
vi quy trình trích xuất và khai thác thông tin trên PDF có cấu trúc bán hình thức,
không phải một hệ thống hỏi đáp tổng quát cho mọi loại tài liệu.

Nhìn chung, đồ án đã chứng minh được ba ý chính: cần tổ chức lại PDF thành dữ
liệu có cấu trúc trước khi hỏi đáp; truy xuất kết hợp cùng siêu dữ liệu nguồn giúp cải
thiện khả năng tìm bằng chứng; và câu trả lời trên PDF cần được đánh giá cùng với
khả năng dẫn chứng, không chỉ bằng độ khớp nội dung cuối cùng.

\section{Hạn chế của đồ án}
\label{section:thesis-limitations}

Mặc dù đạt được các kết quả tích cực trên benchmark chính, đồ án vẫn còn một số
hạn chế rõ ràng.

Khả năng xử lý PDF scan thực tế chưa được đánh giá đầy đủ. Kết quả OCR
trên tập scan tổng hợp cao, nhưng chênh lệch với các tập dữ liệu nhiễu hơn cho thấy
hệ thống vẫn nhạy với chất lượng ảnh, bố cục phức tạp và sai lệch hộp từ. Vì vậy,
đồ án chưa thể khẳng định hiệu quả tương đương trên các tài liệu scan tiếng Việt
thực tế.

Phần tái dựng bảng dù đã cải thiện mạnh ở các độ đo trung gian vẫn chưa
giải quyết trọn vẹn các bảng khó. Exact CSV của Hybrid TATR còn thấp hơn đáng
kể so với Cell F1 và Text Assignment F1, cho thấy hệ thống thường lấy đúng phần
lớn ô và văn bản, nhưng vẫn dễ sai ở các bảng có ô gộp, tiêu đề nhiều hàng hoặc
cấu trúc span phức tạp.

Khả năng tổng hợp nhiều bằng chứng còn hạn chế. Điều này thể hiện rõ trên
QASPER, nơi nhiều câu hỏi yêu cầu kết hợp nhiều đoạn thông tin và tạo đáp án tự
do. Trong các trường hợp như vậy, hệ thống có thể tìm được một phần bằng chứng
đúng nhưng vẫn chưa tổng hợp thành câu trả lời hoàn chỉnh.

Các kết luận mạnh nhất của đồ án hiện chỉ áp dụng cho một số miền dữ liệu
cụ thể. QCDT đại diện cho tài liệu quy định tiếng Việt có lớp văn bản; còn bộ hỏi
đáp trên bảng mở rộng mới bao phủ một tài liệu QCDT.
Vì vậy, các kết quả hiện tại chưa đủ để suy rộng cho mọi loại tài liệu PDF hoặc mọi
bối cảnh nghiệp vụ.

Một giới hạn khác là các độ đo về grounding và hallucination trong đồ án phụ thuộc vào quy
tắc đánh giá đã xác định trước. Chúng rất hữu ích để kiểm tra khả năng dẫn chứng,
nhưng không thay thế hoàn toàn cho đánh giá thủ công về độ đúng, độ đầy đủ và sự
phù hợp ngữ cảnh của câu trả lời.

\section{Hướng phát triển}
\label{section:future-work}

Từ các hạn chế trên, có thể đề xuất một số hướng phát triển tiếp theo.

Hướng thứ nhất là mở rộng đánh giá và cải thiện luồng trích xuất cho PDF scan tiếng Việt
thực tế. Điều này bao gồm xây dựng thêm tập dữ liệu gán nhãn, theo dõi độ tin cậy
OCR theo từng vùng, và kết hợp tốt hơn giữa OCR với định tuyến theo vùng nội
dung.

Hướng thứ hai là tiếp tục cải thiện tái dựng bảng. Các nhóm bảng có ô gộp, tiêu đề
nhiều hàng và cấu trúc không đồng nhất là nơi còn nhiều dư địa nhất. Một hướng đi
khả thi là bổ sung các ràng buộc cấu trúc sau bước nhận diện ô, hoặc dùng biểu diễn
đồ thị để tái dựng quan hệ hàng, cột và span một cách nhất quán hơn.

Hướng thứ ba là tăng cường truy xuất và xếp hạng cho hỏi đáp trên bảng. Kết quả hiện tại cho
thấy nhiều ô đúng đã tồn tại trong tập dữ liệu truy xuất nhưng chưa được đưa lên thứ hạng đủ cao.
Vì vậy, một bộ sắp xếp lại kết quả chuyên biệt cho quan hệ giữa hàng, cột và giá trị ô là một
hướng phát triển phù hợp.

Hướng thứ tư là cải thiện khả năng tổng hợp nhiều bằng chứng ở tầng hỏi đáp. Với
các câu hỏi dài hoặc cần kết hợp nhiều điều kiện, hệ thống cần cơ chế lập kế hoạch
bằng chứng tốt hơn thay vì chỉ dựa vào một lượt truy xuất và sinh câu trả lời ngắn.
Điều này đặc biệt quan trọng nếu muốn mở rộng từ tài liệu quy định sang tài liệu
khoa học hoặc báo cáo dài.

Hướng thứ năm là mở rộng phạm vi đánh giá đầu cuối. Các thí nghiệm tiếp theo nên
bao phủ nhiều tài liệu hơn, nhiều miền hơn và nhiều loại câu hỏi hơn, đặc biệt là
với dữ liệu bảng thực tế và tài liệu tiếng Việt ngoài QCDT. Khi đó, các kết luận về
khả năng tổng quát hóa của hệ thống sẽ vững hơn.

\section{Tổng kết chương}
\label{section:chapter7-summary}

Đồ án đã cho thấy rằng bài toán hỏi đáp trên PDF không thể tách rời khỏi bài toán
trích xuất PDF. Việc tổ chức lại tài liệu thành các block, bảng, ô và siêu dữ liệu nguồn
là điều kiện cần để truy xuất bằng chứng đúng và tạo câu trả lời có thể kiểm chứng.
Trong phạm vi dữ liệu và cấu hình đã đánh giá, hệ thống đạt được kết quả tích cực
trên tập chính và chỉ ra rõ vai trò của từng thành phần trong quy trình. Đồng thời,
đồ án cũng cho thấy những giới hạn còn lại ở xử lý scan, tái dựng bảng phức tạp và
tổng hợp nhiều bằng chứng, từ đó mở ra các hướng phát triển tiếp theo cho bài toán
trích xuất và khai thác thông tin trên tài liệu PDF.

\end{document}
